\section{Discussion and limitations}
The report must identify what was selected, what evaluation information was available during search, and which data support the final score. For a frozen bank, simultaneous intervals provide a valid way to acknowledge selection using the existing search data, but leave point-score optimism intact. A genuinely independent audit estimates the frozen winner directly, while splits can change the winner. Reports should therefore give selection and audit sample sizes, the bank and tie rule, interval construction, and both distinct-question and inference-call costs. A headline score without these details leaves its target and reliability ambiguous.

The controlled experiment shows that selection can alter both numerical accuracy and decisions, with substantial variation across settings. Strong shared errors and true quality separation can make the gap much smaller. A warning based only on the most severe independent equal-quality cell would overstate generality; a favorable high-correlation example would understate the possible failure. The complete grid and its uncertainties expose both cases.

The simulation is intentionally limited. Its mixture dependence does not model all prompt-error patterns. Increasing $K$ in heterogeneous cells also changes available true qualities. Binomial interval guarantees assume iid question draws, fixed scoring rules, and an outcome-independent candidate family. Subject stratification, sampling without replacement, benchmark contamination in pretraining, or feedback-driven candidate creation require separate interpretation. Marginal coverage across searches is not a guarantee conditional on every selected identity or on the fully observed search dataset. The decision endpoint is one conservative nonoverlap rule rather than all scientific comparisons.

A real response cache, when used, supports only the model, quantization, generation settings, prompts, parser, item population, and randomization actually studied. Resampling cached rows reruns selection, not language-model generation; conditional Monte Carlo precision is not uncertainty across future tasks, fresh model samples, or model families. An audit estimate is itself noisy and must not be called known population truth. Claims of population error coverage require a target with known truth or a justified statistical sampling argument, rather than overlap with another estimated score.

The controlled evidence alone does not satisfy the task's required real-LLM experiment. It establishes mechanisms and evaluates reporting protocols under explicit assumptions; evidence from newly generated real responses is necessary before a complete study can assert their magnitude in an actual LLM setting. We claim neither a new discovery of selection bias nor a universally optimal reporting correction.
